\chapter{Related Work}
\label{ch:related_work}

\textbf{Vision-and-Language Navigation.} Early VLN research focused on discrete environments where agents traverse predefined navigation graphs, such as R2R~\cite{anderson2018vision} and RxR~\cite{ku2020room}, before shifting toward Vision-and-Language Navigation in Continuous Environments (VLN-CE)~\cite{krantz2020beyond}, which predicts low-level actions in continuous 3D spaces to better reflect real-world robotic control. Initial VLN-CE approaches relied on cross-modal policies driven by RGB-D and odometry inputs~\cite{krantz2020beyond}, waypoint predictors pretrained in simulators~\cite{an2023etpnav}, or explicit semantic maps and 3D feature fields~\cite{wang2023gridmm,wang2025g3d}.

Recent work has leveraged the generalization capabilities of LLMs and MLLMs to overcome these constraints. Early efforts converted visual observations into text to employ LLMs as zero-shot planners~\cite{zhou2023navgpt,long2024instructnav}, but suffered from severe spatial information loss and high inference latency. NaVid~\cite{zhang2024navid} subsequently established a video-based VLM paradigm that predicts low-level actions end-to-end solely from monocular RGB video without maps, odometry, or depth. Follow-up studies extended this paradigm toward unifying diverse embodied navigation tasks within a single MLLM~\cite{zhang2025uni,cheng2025navila,gao2026octonav,team2026qwen}, while StreamVLN~\cite{wei2026streamvln} achieved low-latency inference over long horizons by combining a sliding-window KV cache with slow-updating memory in a multi-turn dialogue framework. However, because existing MLLM visual encoders are pretrained primarily on 2D image-text pairs, they remain weak at 3D geometric reasoning. To address this limitation, JanusVLN~\cite{zeng2026janusvln} introduced a dual implicit memory architecture that maintains separate initial-window and sliding-window KV caches for an MLLM's 2D semantic encoder and a VGGT-based 3D spatial geometry encoder. Although this substantially improves RGB-only VLN-CE performance, it retains the memory and computational overhead of continuously running an auxiliary geometry encoder at inference time.

\textbf{Geometry Models.} Reconstructing 3D structure from single or multiple RGB images was traditionally formulated as iterative optimization within multi-stage Structure-from-Motion (SfM)~\cite{schnberger2016structure} and Multi-View Stereo (MVS)~\cite{furukawa2015multi} pipelines. Due to their high computational cost and fragile correspondence matching, the field has shifted toward end-to-end learning-based models. DUSt3R~\cite{wang2024dust3r} and its follow-ups~\cite{leroy2024grounding,zhang2025monst3r} marked a turning point by directly regressing 3D pointmaps from image pairs in a feed-forward manner without camera parameters, though scaling to multiple views still required offline global alignment. More recently, multi-view feed-forward geometry foundation models~\cite{yang2025fast3r,wang2025vggt,wang2026permutation,lin2025depth} have enabled single-pass prediction of camera poses and dense 3D geometry from an arbitrary number of images without global alignment, progressing toward permutation-equivariant designs that remove fixed reference-view dependencies and unified representations built on a single plain transformer~\cite{wang2026permutation,lin2025depth}.

Because these multi-view feed-forward models perform bidirectional global attention across all input frames, streaming scenarios require recomputing the entire sequence at each timestep and allow future frame information to leak into current representations. Consequently, recent research has explored online 3D perception for incremental streaming inputs, exemplified by recurrent architectures that update spatial memory and persistent states~\cite{wang20243d,wang2025continuous} and StreamVGGT~\cite{zhuo2026streamvggt}, which replaces global attention with temporal causal attention. Although modern geometry models differ substantially in backbone design, decoding layers, and information flow directionality, how these architectural properties govern VLN distillation efficacy remains largely unexplored.

\textbf{Spatial Knowledge Distillation.} To enhance MLLM spatial perception without dedicated 3D sensors, several studies directly fuse features from pretrained geometry encoders at inference time~\cite{zeng2026janusvln,wu2025spatial,zheng2025learning}, which inevitably incurs external-encoder memory and latency overhead. To eliminate this deployment cost, training-time spatial knowledge distillation based on intermediate representation supervision~\cite{romero2015fitnets} has emerged as an attractive alternative. Aligning frozen geometry teacher features with MLLM visual tokens has expanded from general spatial reasoning~\cite{huang20253drs,huang2026what} and robotic manipulation VLAs~\cite{li2026spatial} to VLN-CE in recent works such as StageVLN~\cite{dao2026stagevln} and GuidedNav~\cite{guo2026guidednav}. However, existing VLN distillation methods predominantly rely on a single VGGT-family teacher and a fixed distillation setup, leaving the impacts of teacher backbone selection, bidirectional-to-causal directional mismatch, and navigation history stride unexamined.
